\section{Consistencia mesoscópica}\label{sec:meso}

La prueba macroscópica usa que el número de líneas del tablero original es de orden $k$.
Cuando quedan pocas filas respecto de $n$, el número de diagonales que tocan el residual
puede ser grande en relación con $k$. En ese régimen conviene mirar cuánta masa llevan:
una diagonal con muchas casillas eliminadas puede tener una masa muy pequeña.

La Sección~\ref{sec:B5T01} impone una cota superior a cada masa diagonal. Bajo esa hipótesis
y las condiciones cuantitativas sobre filas y columnas, las estimaciones se expresan mediante
las masas y las normas de los operadores. El Teorema~\ref{thm:meso} de la
Sección~\ref{sec:B5T05} obtiene así la misma cota de consistencia con constantes
independientes de $n/k$.
